% ============================================================
%  Part VI — Tenses
% ============================================================
\gpart{VI}{Tenses}

% ------------------------------------------------------------
\gsec{The tenses at a glance}

\gintro{German has six tenses but uses them unevenly. Two of them, the Perfekt and the Präsens, carry most of everyday conversation.}

\begin{gtbl}{colspec={X[1,l] X[1.2,l] X[1.3,l]}}
  Tense & Formula & Example \\
  Präsens        & stem + ending            & ich spiele \\
  Perfekt        & haben/sein + Partizip II & ich habe gespielt \\
  Präteritum     & stem + past ending       & ich spielte \\
  Plusquamperfekt & hatte/war + Partizip II & ich hatte gespielt \\
  Futur I        & werden + Infinitiv       & ich werde spielen \\
  Futur II       & werden + Part.\,II + haben/sein & ich werde gespielt haben \\
\end{gtbl}

\tcap{Which one do I actually use?}

\begin{flattbl}{colspec={X[1,l] X[2,l]}}
  Speaking about the past & \textbf{Perfekt} — nearly always \\
  Writing about the past  & \textbf{Präteritum} — narrative, reports \\
  Either way              & \textit{sein}, \textit{haben} and the modals prefer the Präteritum even in speech \\
  The future              & \textbf{Präsens} + a time word, unless you need emphasis \\
\end{flattbl}

\begin{gnote}
\textit{Ich gehe morgen ins Kino} is the normal way to talk about tomorrow.
Reserve \textit{werden} for genuine predictions, promises and assumptions:
\textit{Es wird regnen} — \textit{Sie wird schon zu Hause sein}.
\end{gnote}

% ------------------------------------------------------------
\gsec{Perfekt}

\gintro{The Perfekt is the ordinary past tense of spoken German. It is built from an auxiliary in second position and a past participle at the end of the clause.}

The everyday past tense. Auxiliary in second position, participle at the end.

\begin{gtbl}{colspec={X[0.95,l] X[2.05,l]}}
  Element & Example \\
  Auxiliary + Partizip II & Ich \textbf{habe} das Buch \textbf{gelesen}. \\
  With sein               & Ich \textbf{bin} nach Berlin \textbf{gefahren}. \\
  Subordinate clause      & \dots, weil ich das Buch gelesen \textbf{habe}. \\
\end{gtbl}

\tcap{Forming the Partizip II}

\begin{gtbl}{colspec={X[0.85,l] X[1.15,l] X[1,l]}}
  Verb type & Pattern & Example \\
  Weak            & ge- + stem + \en{-t}  & gespielt, gemacht \\
  Weak, -d/-t stem & ge- + stem + \en{-et} & gearbeitet, geredet \\
  Strong          & ge- + stem + \en{-en} & gesehen, gefahren \\
  Mixed           & ge- + new stem + \en{-t} & gebracht, gedacht \\
  Separable       & prefix + \en{ge} + rest & auf\textbf{ge}standen \\
  Inseparable     & \textbf{no ge-}       & verstanden, bezahlt \\
  -ieren          & \textbf{no ge-}       & studiert, telefoniert \\
\end{gtbl}

\tcap{haben or sein?}

\begin{gtbl}{colspec={X[0.75,l] X[1.15,l] X[1.1,l]}}
  Auxiliary & When & Examples \\
  sein  & change of place  & gehen, fahren, kommen, fliegen \\
  sein  & change of state  & aufstehen, einschlafen, sterben, wachsen \\
  sein  & three specials   & sein, bleiben, werden \\
  haben & everything else  & essen, lesen, machen, haben \\
  haben & all reflexives   & sich freuen, sich waschen \\
\end{gtbl}

\begin{gnote}
The sein-verbs are all intransitive — they never take a direct object. If the
verb has an accusative object, the auxiliary is \textit{haben}. Compare
\textit{Ich \textbf{bin} gefahren} with \textit{Ich \textbf{habe} das Auto
gefahren}.
\end{gnote}

% ------------------------------------------------------------
\gsec{Präteritum}

\gintro{The Präteritum is the one-word past tense. It belongs to writing and narrative, with the exception of \textit{sein}, \textit{haben} and the modals, which use it in speech too.}

\tcap{Weak verbs — insert -t- then add the ending}

\begin{gtbl}{colspec={X[0.9,l] X[0.8,c] X[1,c] X[1.1,c]}}
  Person & Ending & spielen & arbeiten \\
  ich       & \en{-te}   & spielte   & arbeitete   \\
  du        & \en{-test} & spieltest & arbeitetest \\
  er/sie/es & \en{-te}   & spielte   & arbeitete   \\
  wir       & \en{-ten}  & spielten  & arbeiteten  \\
  ihr       & \en{-tet}  & spieltet  & arbeitetet  \\
  sie/Sie   & \en{-ten}  & spielten  & arbeiteten  \\
\end{gtbl}

\tcap{Strong verbs — new stem, and no ending in 1st/3rd singular}

\begin{gtbl}{colspec={X[0.9,l] X[0.8,c] X[1,c] X[1,c]}}
  Person & Ending & gehen \eng{ging} & sehen \eng{sah} \\
  ich       & \en{--}   & ging   & sah    \\
  du        & \en{-st}  & gingst & sahst  \\
  er/sie/es & \en{--}   & ging   & sah    \\
  wir       & \en{-en}  & gingen & sahen  \\
  ihr       & \en{-t}   & gingt  & saht   \\
  sie/Sie   & \en{-en}  & gingen & sahen  \\
\end{gtbl}

\begin{gnote}
Mixed verbs take the strong stem change \emph{and} the weak endings:
\textit{bringen} $\rightarrow$ \textit{brachte}, \textit{denken} $\rightarrow$
\textit{dachte}, \textit{wissen} $\rightarrow$ \textit{wusste},
\textit{kennen} $\rightarrow$ \textit{kannte}.
\end{gnote}

% ------------------------------------------------------------
\gsec{Plusquamperfekt and Futur}

\gintro{The pluperfect places one past event before another. The future is formed with \textit{werden}, though the present tense usually does the same job.}

\tcap{Plusquamperfekt — the past before the past}

\begin{gtbl}{colspec={X[0.9,l] X[2.1,l]}}
  Formula & \textbf{hatte} / \textbf{war} + Partizip II \\
  Example & Ich \textbf{hatte} das Buch schon \textbf{gelesen}. \\
  Example & Er \textbf{war} bereits \textbf{gegangen}. \\
  Trigger & almost always paired with \textit{nachdem} \\
\end{gtbl}

\begin{gnote}
\textit{Nachdem} forces a tense step down: \textit{Nachdem ich gegessen
\textbf{hatte}, ging ich ins Bett.} Past-before-past in the subordinate clause,
simple past in the main clause.
\end{gnote}

\tcap{Futur I — werden + Infinitiv}

\begin{gtbl}{colspec={X[0.9,l] X[1,c] X[1.4,l]}}
  Person & werden & Example \\
  ich       & werde  & Ich \textbf{werde} morgen \textbf{kommen}. \\
  du        & wirst  & Du \textbf{wirst} es \textbf{schaffen}. \\
  er/sie/es & wird   & Es \textbf{wird} \textbf{regnen}. \\
  wir       & werden & Wir \textbf{werden} \textbf{sehen}. \\
  ihr       & werdet & Ihr \textbf{werdet} \textbf{lachen}. \\
  sie/Sie   & werden & Sie \textbf{werden} \textbf{warten}. \\
\end{gtbl}

\tcap{Futur II — will have done}

\begin{flattbl}{colspec={X[0.85,l] X[2.15,l]}}
  Formula & werden + Partizip II + haben/sein \\
  Example & Bis Montag \textbf{werde} ich es \textbf{fertig gemacht haben}. \\
  Also    & often expresses an assumption: \textit{Er wird es vergessen haben.} \\
\end{flattbl}

\begin{gnote}
\textit{werden} plus \textit{wohl}, \textit{schon} or \textit{sicher} usually
means \emph{probably}, not \emph{will}: \textit{Sie wird wohl krank sein} —
\emph{she is probably ill}.
\end{gnote}
